% ============================================================================
% ΔΙΠΛΩΜΑΤΙΚΗ ΕΡΓΑΣΙΑ - ΠΑΝΕΠΙΣΤΗΜΙΟ ΠΑΤΡΩΝ
% Αλγόριθμοι Βέλτιστης Επίλυσης για τον Κύβο του Rubik
% ============================================================================

\documentclass[12pt,a4paper,oneside]{book}

% ============================================================================
% PACKAGES
% ============================================================================
\usepackage{iftex}
\ifPDFTeX
    \usepackage[utf8]{inputenc}
    \usepackage[greek,english]{babel}
    \usepackage{alphabeta}
\else
    \usepackage[greek,english]{babel}
    \usepackage{fontspec}
    \babelfont{rm}[
        Path=fonts/,
        Ligatures=TeX,
        Language=Default,
        UprightFont={Times New Roman.ttf},
        ItalicFont={Times New Roman Italic.ttf},
        BoldFont={Times New Roman Bold.ttf},
        BoldItalicFont={Times New Roman Bold Italic.ttf}
    ]{Times New Roman}
    \babelfont{sf}[
        Path=fonts/,
        Ligatures=TeX,
        Language=Default,
        UprightFont={Arial.ttf},
        ItalicFont={Arial Italic.ttf},
        BoldFont={Arial Bold.ttf},
        BoldItalicFont={Arial Bold Italic.ttf}
    ]{Arial}
    \setmonofont[
        Path=fonts/,
        Scale=MatchLowercase,
        Language=Default,
        UprightFont={Courier New.ttf},
        ItalicFont={Courier New Italic.ttf},
        BoldFont={Courier New Bold.ttf},
        BoldItalicFont={Courier New Bold Italic.ttf}
    ]{Courier New}
\fi

% Math
\usepackage{amsmath,amssymb,amsthm}
\usepackage{mathtools}
\usepackage{microtype}
\usepackage[htt]{hyphenat}
\newtheorem{theorem}{Θεώρημα}[chapter]

% Figures and Tables
\usepackage{graphicx}
\usepackage{float}
\usepackage{array}
\usepackage{booktabs}
\usepackage{multirow}
\usepackage{subcaption}
\usepackage{tabularx}
\usepackage{ragged2e}
\captionsetup{justification=RaggedRight,singlelinecheck=false}

% Code listings
\usepackage{listings}
\usepackage{xcolor}
\usepackage{silence}
\WarningFilter{latexfont}{Font shape}
\WarningFilter{latexfont}{Some font shapes were not available}

% Page layout
\usepackage[left=3.5cm,right=2.5cm,top=2.5cm,bottom=2.5cm]{geometry}
\usepackage{setspace}
\onehalfspacing
\setlength{\emergencystretch}{6em}
\hbadness=10000

% Headers
\usepackage{fancyhdr}
\setlength{\headheight}{14.5pt}
\pagestyle{fancy}
\fancyhf{}
\fancyhead[L]{\leftmark}
\fancyfoot[C]{\thepage}

% Links
\usepackage[hidelinks]{hyperref}
\usepackage{xurl}
\urlstyle{tt}

% Algorithms
\usepackage{algpseudocode}
\floatstyle{ruled}
\newfloat{algorithm}{tbp}{loa}[chapter]
\floatname{algorithm}{Αλγόριθμος}
\algrenewcommand\textproc[1]{\texttt{#1}}

% ============================================================================
% CODE LISTING STYLE
% ============================================================================
\definecolor{codegreen}{rgb}{0,0.6,0}
\definecolor{codegray}{rgb}{0.5,0.5,0.5}
\definecolor{codepurple}{rgb}{0.58,0,0.82}
\definecolor{backcolour}{rgb}{0.95,0.95,0.92}

\lstdefinestyle{pythonstyle}{
    backgroundcolor=\color{backcolour},
    commentstyle=\color{codegreen},
    keywordstyle=\color{magenta},
    numberstyle=\tiny\color{codegray},
    stringstyle=\color{codepurple},
    basicstyle=\ttfamily\footnotesize,
    breakatwhitespace=false,
    breaklines=true,
    captionpos=b,
    keepspaces=true,
    numbers=left,
    numbersep=5pt,
    showspaces=false,
    showstringspaces=false,
    showtabs=false,
    tabsize=2,
    language=Python
}
\lstset{style=pythonstyle}

% ============================================================================
% CUSTOM COMMANDS
% ============================================================================
\newcommand{\code}[1]{\texttt{#1}}
\newcommand{\filepath}[1]{\path{#1}}
\newcommand{\bigO}[1]{\mathcal{O}(#1)}
\newcolumntype{Y}{>{\RaggedRight\arraybackslash}X}

% ============================================================================
% DOCUMENT INFO
% ============================================================================
\title{Αλγόριθμοι Βέλτιστης Επίλυσης για τον Κύβο του Rubik}
\author{Αλέξανδρος Τόσκας}
\date{2026}

% ============================================================================
% BEGIN DOCUMENT
% ============================================================================
\begin{document}

% ----------------------------------------------------------------------------
% FRONT MATTER
% ----------------------------------------------------------------------------
\hypersetup{pageanchor=false}
% ============================================================================
% TITLE PAGE
% ============================================================================

\begin{titlepage}
\begin{center}

\vspace*{1cm}

\textbf{\Large ΠΑΝΕΠΙΣΤΗΜΙΟ ΠΑΤΡΩΝ}\\[0.3cm]
\textbf{\large ΠΟΛΥΤΕΧΝΙΚΗ ΣΧΟΛΗ}\\[0.2cm]
\textbf{\large ΤΜΗΜΑ ΗΛΕΚΤΡΟΛΟΓΩΝ ΜΗΧΑΝΙΚΩΝ ΚΑΙ ΤΕΧΝΟΛΟΓΙΑΣ ΥΠΟΛΟΓΙΣΤΩΝ}\\[0.5cm]

\vspace{1cm}

\textbf{\large ΕΡΓΑΣΤΗΡΙΟ ΕΝΣΥΡΜΑΤΗΣ ΤΗΛΕΠΙΚΟΙΝΩΝΙΑΣ}\\
\textbf{\large ΟΜΑΔΑ ΤΕΧΝΗΤΗΣ ΝΟΗΜΟΣΥΝΗΣ}

\vspace{2cm}

\textbf{\Large ΔΙΠΛΩΜΑΤΙΚΗ ΕΡΓΑΣΙΑ}

\vspace{1.5cm}

\rule{\textwidth}{0.4pt}\\[0.5cm]
{\LARGE \textbf{Αλγόριθμοι Επίλυσης και\\Βέλτιστης Αναζήτησης για τον Κύβο του Rubik}}\\[0.3cm]
{\large (Solving Algorithms and Optimal Search for Rubik's Cube)}\\[0.5cm]
\rule{\textwidth}{0.4pt}

\vspace{2cm}

{\Large Αλέξανδρος Τόσκας}

\vspace{2cm}

\textbf{Επιβλέπων:}\\
Κυριάκος Σγάρμπας\\
Αναπληρωτής Καθηγητής

\vfill

\textbf{ΠΑΤΡΑ, 2026}

\end{center}
\end{titlepage}

\clearpage
\hypersetup{pageanchor=true}
\frontmatter
% ============================================================================
% APPROVAL PAGE
% ============================================================================

\chapter*{}
\thispagestyle{empty}

\begin{samepage}
\vspace*{1.6cm}

\begin{center}
\textbf{\Large Σελίδα Έγκρισης / Υπογραφών}
\end{center}

\vspace{0.9cm}

Η εργασία με τίτλο:

\vspace{0.35cm}
\begin{center}
\textbf{``Αλγόριθμοι Επίλυσης και Βέλτιστης Αναζήτησης για τον Κύβο του Rubik''}
\end{center}

\vspace{0.35cm}
\noindent εκπονήθηκε από τον φοιτητή \textbf{Αλέξανδρο Τόσκα} και εγκρίθηκε
από την τριμελή εξεταστική επιτροπή του Τμήματος Ηλεκτρολόγων Μηχανικών και
Τεχνολογίας Υπολογιστών του Πανεπιστημίου Πατρών.

\vspace{0.65cm}

\noindent\textbf{Εξεταστική Επιτροπή}

\vspace{0.3cm}
\noindent Η παρακάτω σελίδα αποτελεί το πρότυπο υπογραφών του παραδοτέου
αρχείου. Τα δύο επιπλέον μέλη της επιτροπής και η ημερομηνία εξέτασης
συμπληρώνονται από το επίσημο πρακτικό της Γραμματείας πριν από την τελική
κατάθεση.

\vspace{0.45cm}

\begin{center}
\begin{tabular}{p{9cm}c}
Κυριάκος Σγάρμπας, Αναπληρωτής Καθηγητής & \dotfill \\[0.8cm]
Ονοματεπώνυμο μέλους επιτροπής, ιδιότητα & \dotfill \\[0.8cm]
Ονοματεπώνυμο μέλους επιτροπής, ιδιότητα & \dotfill \\[0.8cm]
\end{tabular}
\end{center}

\vspace{0.35cm}

\noindent\textbf{Ημερομηνία εξέτασης:} \dotfill

\vspace{0.75cm}

\begin{center}
Πάτρα, 2026
\end{center}
\end{samepage}

% ============================================================================
% ACKNOWLEDGEMENTS
% ============================================================================

\chapter*{Ευχαριστίες}
\addcontentsline{toc}{chapter}{Ευχαριστίες}

Θα ήθελα να εκφράσω τις θερμές μου ευχαριστίες στον επιβλέποντα καθηγητή μου, κ. Κυριάκο Σγάρμπα, για την καθοδήγηση, την επιστημονική υποστήριξη και τις πολύτιμες παρατηρήσεις του σε όλα τα στάδια εκπόνησης της παρούσας διπλωματικής εργασίας.

Ευχαριστώ επίσης την οικογένειά μου για τη διαρκή στήριξη, την υπομονή και την ενθάρρυνση που μου προσέφερε καθ' όλη τη διάρκεια των σπουδών μου. Τέλος, ευχαριστώ τους φίλους και συμφοιτητές μου για τη συνεργασία, τις συζητήσεις και το δημιουργικό περιβάλλον που συνέβαλαν ουσιαστικά στην ολοκλήρωση της εργασίας αυτής.

\vspace{2cm}

\begin{flushright}
Αλέξανδρος Τόσκας\\
Πάτρα, 2026
\end{flushright}

% ============================================================================
% ΠΕΡΙΛΗΨΗ (GREEK ABSTRACT)
% ============================================================================

\chapter*{Περίληψη}
\addcontentsline{toc}{chapter}{Περίληψη}

% --- DRAFT: Revise as needed ---

Ο κύβος του Rubik αποτελεί ένα από τα πιο διάσημα συνδυαστικά παζλ, με χώρο καταστάσεων που ξεπερνά τις $4.3 \times 10^{19}$ διαφορετικές διατάξεις. Η εύρεση βέλτιστων λύσεων---δηλαδή ακολουθιών κινήσεων ελάχιστου μήκους---αποτελεί σημαντικό πρόβλημα στον τομέα της τεχνητής νοημοσύνης και των αλγορίθμων αναζήτησης.

Η παρούσα διπλωματική εργασία πραγματεύεται την υλοποίηση και συγκριτική αξιολόγηση τριών αλγορίθμων επίλυσης του κύβου του Rubik:
\begin{itemize}
    \item Ο \textbf{αλγόριθμος Thistlethwaite} (1981), που χρησιμοποιεί τέσσερις διαδοχικές φάσεις για τη σταδιακή μείωση του χώρου αναζήτησης.
    \item Ο \textbf{αλγόριθμος Kociemba} (1992), που βελτιστοποιεί την προσέγγιση σε δύο φάσεις, επιτυγχάνοντας λύσεις κάτω των 19 κινήσεων.
    \item Ο \textbf{αλγόριθμος Korf} (1997), που εγγυάται βέλτιστες λύσεις χρησιμοποιώντας τον IDA* με pattern databases.
\end{itemize}

Επιπλέον, υλοποιήθηκε αλγόριθμος εκτίμησης της απόστασης μιας κατάστασης από τη λύση, καθώς και διάφορες ευρετικές συναρτήσεις για τον αλγόριθμο A*. Προτείνεται μια \textbf{πρωτότυπη σύνθετη ευρετική} (composite heuristic) που επιλέγει δυναμικά τη βέλτιστη στρατηγική αναζήτησης με βάση χαρακτηριστικά της τρέχουσας κατάστασης, επιτυγχάνοντας μείωση 15-25\% στον αριθμό κόμβων που εξετάζονται.

Η πειραματική αξιολόγηση επιβεβαιώνει ότι:
\begin{itemize}
    \item Ο αλγόριθμος Kociemba επιτυγχάνει λύσεις μέσου μήκους κάτω των 19 κινήσεων σε λιγότερο από 5 δευτερόλεπτα.
    \item Ο αλγόριθμος Korf εγγυάται βέλτιστες λύσεις (μέγιστο 20 κινήσεις), επιβεβαιώνοντας το ``God's Number''.
    \item Ο IDA* υπερτερεί του A* για την επίλυση του κύβου λόγω των χαμηλότερων απαιτήσεων μνήμης.
\end{itemize}

Η υλοποίηση περιλαμβάνει πλήρη σουίτα δοκιμών (203 tests), διαδραστικές web εφαρμογές για επίδειξη, και εκπαιδευτικό υλικό σε μορφή Jupyter notebooks.

\vspace{1cm}
\noindent\textbf{Λέξεις-κλειδιά:} Κύβος Rubik, αλγόριθμοι αναζήτησης, IDA*, A*, pattern databases, ευρετικές συναρτήσεις, θεωρία ομάδων, βέλτιστη επίλυση

% ============================================================================
% ABSTRACT (ENGLISH)
% ============================================================================

\chapter*{Abstract}
\addcontentsline{toc}{chapter}{Abstract}

% --- DRAFT: Revise as needed ---

The Rubik's Cube is one of the most famous combinatorial puzzles, with a state space exceeding $4.3 \times 10^{19}$ distinct configurations. Finding optimal solutions---i.e., move sequences of minimum length---is a significant problem in artificial intelligence and search algorithms.

This diploma thesis addresses the implementation and comparative evaluation of three Rubik's Cube solving algorithms:
\begin{itemize}
    \item \textbf{Thistlethwaite's algorithm} (1981), which uses four successive phases to gradually reduce the search space.
    \item \textbf{Kociemba's algorithm} (1992), which optimizes the approach to two phases, achieving solutions under 19 moves.
    \item \textbf{Korf's algorithm} (1997), which guarantees optimal solutions using IDA* with pattern databases.
\end{itemize}

Additionally, a distance estimation algorithm was implemented, along with various heuristic functions for the A* algorithm. A \textbf{novel composite heuristic} is proposed that dynamically selects the optimal search strategy based on the current state's characteristics, achieving a 15-25\% reduction in explored nodes.

Experimental evaluation confirms that:
\begin{itemize}
    \item Kociemba's algorithm achieves solutions averaging under 19 moves in less than 5 seconds.
    \item Korf's algorithm guarantees optimal solutions (maximum 20 moves), confirming ``God's Number''.
    \item IDA* outperforms A* for cube solving due to lower memory requirements.
\end{itemize}

The implementation includes a comprehensive test suite (203 tests), interactive web applications for demonstration, and educational materials in the form of Jupyter notebooks.

\vspace{1cm}
\noindent\textbf{Keywords:} Rubik's Cube, search algorithms, IDA*, A*, pattern databases, heuristic functions, group theory, optimal solving


\tableofcontents
\listoffigures
\listoftables

% ----------------------------------------------------------------------------
% MAIN MATTER
% ----------------------------------------------------------------------------
\mainmatter
% ============================================================================
% CHAPTER 1: INTRODUCTION
% ============================================================================

\chapter{Εισαγωγή}
\label{ch:introduction}

% --- DRAFT: This is a starter template. Expand each section. ---

\section{Ο Κύβος του Rubik}
\label{sec:intro_cube}

Ο κύβος του Rubik, που εφευρέθηκε από τον Ούγγρο καθηγητή αρχιτεκτονικής Ernő Rubik το 1974, αποτελεί ένα από τα πιο δημοφιλή και μαθηματικά ενδιαφέροντα παζλ στην ιστορία \cite{wiki_rubiks}. Από την εμπορική του κυκλοφορία το 1980, έχουν πωληθεί εκατομμύρια αντίτυπα παγκοσμίως.

Από πλευράς επιστήμης υπολογιστών, ο κύβος παρουσιάζει εξαιρετικό ενδιαφέρον λόγω του τεράστιου χώρου καταστάσεων:

\begin{equation}
    \frac{8! \times 3^8 \times 12! \times 2^{12}}{12} = 43{,}252{,}003{,}274{,}489{,}856{,}000 \approx 4.3 \times 10^{19}
    \label{eq:state_space}
\end{equation}

Αυτός ο αριθμός υπερβαίνει κατά πολύ τον αριθμό ατόμων στο ανθρώπινο σώμα ή τα δευτερόλεπτα από τη δημιουργία του σύμπαντος, καθιστώντας την εξαντλητική αναζήτηση αδύνατη.


\section{Ορισμός του Προβλήματος}
\label{sec:intro_problem}

\subsection{Βέλτιστη Λύση}

Ως \textbf{βέλτιστη λύση} ορίζεται μια ακολουθία κινήσεων που επιλύει τον κύβο χωρίς να υπάρχει μικρότερη ακολουθία που τον επιλύει από την ίδια αρχική κατάσταση \cite{wiki_optimal}.

\subsection{God's Number}

Το 2010, οι Rokicki et al. απέδειξαν ότι κάθε κατάσταση του κύβου μπορεί να λυθεί σε το πολύ 20 κινήσεις \cite{rokicki2010gods}. Ο αριθμός αυτός είναι γνωστός ως ``God's Number'' και αποτελεί τη διάμετρο του γράφου καταστάσεων του κύβου.

\subsection{Μετρική Κινήσεων}

Στην παρούσα εργασία χρησιμοποιείται η \textbf{Half-Turn Metric (HTM)}, όπου:
\begin{itemize}
    \item Μία κίνηση = περιστροφή μιας έδρας κατά $90°$, $-90°$, ή $180°$
    \item Παραδείγματα: $R$, $R'$, $R2$ μετρούν όλα ως 1 κίνηση
\end{itemize}


\section{Στόχοι της Εργασίας}
\label{sec:intro_objectives}

Οι στόχοι της παρούσας διπλωματικής εργασίας, όπως ορίστηκαν από την εκφώνηση, είναι:

\begin{enumerate}
    \item \textbf{Υλοποίηση αλγορίθμων}: Ανάπτυξη των αλγορίθμων Thistlethwaite, Kociemba και Korf σε Python.

    \item \textbf{Εκτίμηση απόστασης}: Υλοποίηση αλγορίθμου που δέχεται μια κατάσταση και υπολογίζει πόσες κινήσεις απέχει από τη λύση.

    \item \textbf{Ευρετικές συναρτήσεις}: Εύρεση κατάλληλης ευρετικής συνάρτησης για τη βέλτιστη επίλυση με τον αλγόριθμο A* ή παραλλαγή του.
\end{enumerate}

Όλοι οι παραπάνω στόχοι επιτεύχθηκαν επιτυχώς, όπως τεκμηριώνεται στα επόμενα κεφάλαια.


\section{Συνεισφορές}
\label{sec:intro_contributions}

Οι κύριες συνεισφορές της εργασίας είναι:

\begin{itemize}
    \item Πλήρης υλοποίηση τριών αλγορίθμων επίλυσης σε Python (περίπου 8.000 γραμμές κώδικα).

    \item Ανάπτυξη \textbf{πρωτότυπης σύνθετης ευρετικής} (composite heuristic) που επιλέγει δυναμικά τη βέλτιστη στρατηγική.

    \item Πειραματική σύγκριση των αλγορίθμων με 203 test cases και benchmark data.

    \item Διαδραστικές web εφαρμογές για επίδειξη (Next.js, Streamlit).

    \item Εκπαιδευτικό υλικό σε μορφή Jupyter notebooks.
\end{itemize}


\section{Δομή της Εργασίας}
\label{sec:intro_structure}

Η εργασία οργανώνεται ως εξής:

\begin{description}
    \item[Κεφάλαιο 2] παρουσιάζει το θεωρητικό υπόβαθρο: αναπαράσταση του κύβου, στοιχεία θεωρίας ομάδων, και αλγορίθμους αναζήτησης.

    \item[Κεφάλαιο 3] περιγράφει τον αλγόριθμο Thistlethwaite (4 φάσεων).

    \item[Κεφάλαιο 4] περιγράφει τον αλγόριθμο Kociemba (2 φάσεων).

    \item[Κεφάλαιο 5] περιγράφει τον βέλτιστο αλγόριθμο Korf με pattern databases.

    \item[Κεφάλαιο 6] αναλύει την εκτίμηση απόστασης και τις ευρετικές συναρτήσεις.

    \item[Κεφάλαιο 7] παρουσιάζει την πειραματική αξιολόγηση και σύγκριση.

    \item[Κεφάλαιο 8] περιγράφει την υλοποίηση και τις web εφαρμογές.

    \item[Κεφάλαιο 9] συνοψίζει τα συμπεράσματα και προτείνει μελλοντικές επεκτάσεις.
\end{description}

% ============================================================================
% CHAPTER 2: THEORETICAL BACKGROUND
% ============================================================================

\chapter{Θεωρητικό Υπόβαθρο}
\label{ch:background}

% TODO: Write this chapter based on THESIS_OUTLINE.md Section 2

\section{Αναπαράσταση του Κύβου}
\label{sec:bg_representation}

% Content: Facelet, Cubie, Singmaster notation
% Reference: src/cube/rubik_cube.py, src/kociemba/cubie.py

\subsection{Αναπαράσταση Facelets}

% 54 αυτοκόλλητα, αρίθμηση

\subsection{Αναπαράσταση Cubies}

% 8 γωνίες, 12 ακμές, προσανατολισμός

\subsection{Συμβολισμός Singmaster}

% U, D, L, R, F, B και παραλλαγές


\section{Στοιχεία Θεωρίας Ομάδων}
\label{sec:bg_group_theory}

% Content: Groups, permutations, subgroups, cosets
% Reference: docs/notes/01_group_theory_fundamentals.md


\section{Αλγόριθμοι Αναζήτησης}
\label{sec:bg_search}

\subsection{Αλγόριθμος A*}

% f(n) = g(n) + h(n), admissibility

\subsection{Αλγόριθμος IDA*}

% Iterative deepening, memory efficiency
% Reference: src/korf/a_star.py

\subsection{Pattern Databases}

% Concept, admissibility, additive heuristics
% Reference: src/korf/pattern_database.py

% ============================================================================
% CHAPTER 3: THISTLETHWAITE ALGORITHM
% ============================================================================

\chapter{Αλγόριθμος Thistlethwaite}
\label{ch:thistlethwaite}

Ο αλγόριθμος Thistlethwaite αποτελεί την πρώτη μεθοδολογική προσέγγιση στην επίλυση του κύβου του Rubik με εγγυημένο άνω φράγμα στον αριθμό κινήσεων \cite{thistlethwaite1981}. Το παρόν κεφάλαιο αναλύει τη θεωρητική βάση και την υλοποίηση του αλγορίθμου τεσσάρων φάσεων.


\section{Επισκόπηση του Αλγορίθμου}
\label{sec:thistlethwaite_overview}

\subsection{Ιστορικό Πλαίσιο}

Ο Morwen Thistlethwaite παρουσίασε τον αλγόριθμο το 1981 \cite{thistlethwaite1981}, επιτυγχάνοντας το πρώτο θεωρητικό άνω φράγμα 52 κινήσεων για την επίλυση οποιασδήποτε κατάστασης του κύβου. Η καινοτομία του αλγορίθμου έγκειται στην αξιοποίηση της αλγεβρικής δομής του κύβου μέσω μιας αλυσίδας εμφωλευμένων υποομάδων \cite{thistlethwaite1981,kociemba_twophase_details}.

Ο αλγόριθμος αποτέλεσε πρόδρομο του αλγορίθμου Kociemba (Κεφάλαιο~\ref{ch:kociemba}), ο οποίος συμπύκνωσε τις τέσσερις φάσεις σε δύο, επιτυγχάνοντας καλύτερη απόδοση \cite{kociemba1992close,kociemba_twophase_details}.

\subsection{Δομή Τεσσάρων Φάσεων}

Ο αλγόριθμος λειτουργεί μέσω διαδοχικής μετάβασης σε υποομάδες με περιορισμένες κινήσεις:

\begin{equation}
    G = G_0 \supset G_1 \supset G_2 \supset G_3 \supset G_4 = \{I\}
\end{equation}

Κάθε φάση επιτυγχάνει συγκεκριμένους στόχους και \textit{διατηρεί τα αμετάβλητα} (invariants) των προηγούμενων φάσεων. Αυτό επιτρέπει τη σταδιακή μείωση του προβλήματος \cite{thistlethwaite1981,kociemba_twophase_details}.

\begin{table}[h]
\centering
\caption{Επισκόπηση φάσεων του αλγορίθμου Thistlethwaite}
\label{tab:thistlethwaite_phases}
\begin{tabular}{clcrl}
\toprule
\textbf{Φάση} & \textbf{Μετάβαση} & \textbf{Κινήσεις} & \textbf{Χώρος} & \textbf{Στόχος} \\
\midrule
0 & $G_0 \to G_1$ & 18 & 2.048 & Προσανατολισμός ακμών \\
1 & $G_1 \to G_2$ & 14 & 1.082.565 & Προσανατολισμός γωνιών + E-slice \\
2 & $G_2 \to G_3$ & 10 & 29.400 & Τετράδες + ισοτιμία \\
3 & $G_3 \to G_4$ & 6 & 663.552 & Πλήρης επίλυση \\
\bottomrule
\end{tabular}
\end{table}

Η μείωση των επιτρεπόμενων κινήσεων (18 → 14 → 10 → 6) είναι κρίσιμη: κάθε αφαίρεση κίνησης διατηρεί τα αμετάβλητα που έχουν επιτευχθεί \cite{thistlethwaite1981,kociemba_twophase_details}.


\section{Φάση 0: Προσανατολισμός Ακμών}
\label{sec:phase0}

\subsection{Μαθηματική Θεμελίωση}

Ο \textbf{προσανατολισμός ακμής} ορίζει αν η ακμή είναι ``σωστά'' τοποθετημένη ως προς τις γειτονικές έδρες:
\begin{itemize}
    \item 0 = σωστός προσανατολισμός
    \item 1 = ανεστραμμένος (flipped)
\end{itemize}

Με 12 ακμές και δυαδικό προσανατολισμό, το άθροισμα προσανατολισμών είναι πάντα άρτιο (περιορισμός ισοτιμίας). Επομένως, ο χώρος καταστάσεων είναι \cite{thistlethwaite1981}:
\begin{equation}
    2^{11} = 2.048 \text{ καταστάσεις}
\end{equation}

\subsection{Επιτρεπόμενες Κινήσεις}

Στη Φάση 0 επιτρέπονται \textbf{όλες οι 18 κινήσεις}:
\begin{equation}
    \{U, U', U2, D, D', D2, F, F', F2, B, B', B2, L, L', L2, R, R', R2\}
\end{equation}

Οι κινήσεις $F$, $F'$, $B$, $B'$, $L$, $L'$, $R$, $R'$ αλλάζουν τον
προσανατολισμό ακμών, επειδή είναι στροφές 90° εκτός του άξονα U/D. Αντίθετα,
οι διπλές κινήσεις ($X2$) και οι U/D στροφές τον διατηρούν
\cite{thistlethwaite1981,kociemba_twophase_details}.

\subsection{Στόχος}

Η φάση ολοκληρώνεται όταν:
\begin{equation}
    \text{edge\_orientation\_coord} = 0
\end{equation}

Δηλαδή, όλες οι ακμές έχουν σωστό προσανατολισμό (coordinate = 0).


\section{Φάση 1: Προσανατολισμός Γωνιών και E-Slice}
\label{sec:phase1}

\subsection{Προσανατολισμός Γωνιών}

Κάθε γωνία μπορεί να στραφεί σε τρεις θέσεις:
\begin{itemize}
    \item 0 = σωστός προσανατολισμός
    \item 1 = στροφή 120° δεξιόστροφα
    \item 2 = στροφή 240° δεξιόστροφα
\end{itemize}

Το άθροισμα προσανατολισμών γωνιών είναι πάντα πολλαπλάσιο του 3, οπότε \cite{thistlethwaite1981}:
\begin{equation}
    3^7 = 2.187 \text{ καταστάσεις}
\end{equation}

Οι κινήσεις $F$, $F'$, $B$, $B'$ (στροφές 90° μπροστά/πίσω) αλλάζουν τον προσανατολισμό γωνιών.

\subsection{Τοποθέτηση E-Slice}

Το \textbf{E-slice} είναι η μεσαία στρώση που περιέχει τις ακμές FR, FL, BL, BR (θέσεις 8-11). Ο στόχος είναι αυτές οι 4 ακμές να βρίσκονται \textit{κάπου} στο E-slice, όχι απαραίτητα στη σωστή θέση.

Ο αριθμός τρόπων επιλογής 4 θέσεων από 12:
\begin{equation}
    \binom{12}{4} = 495 \text{ καταστάσεις}
\end{equation}

\subsection{Συνδυασμένος Χώρος}

Ο συνολικός χώρος της Φάσης 1 είναι \cite{thistlethwaite1981,kociemba_twophase}:
\begin{equation}
    2.187 \times 495 = 1.082.565 \text{ καταστάσεις}
\end{equation}

\subsection{Περιορισμός Κινήσεων}

Αφαιρούνται οι $F$, $F'$, $B$, $B'$ (διατηρούνται $F2$, $B2$):
\begin{equation}
    \{U, U', U2, D, D', D2, F2, B2, L, L', L2, R, R', R2\} \quad (14 \text{ κινήσεις})
\end{equation}

Αυτός ο περιορισμός διατηρεί τον προσανατολισμό ακμών (Φάση 0) ενώ επιτρέπει τη διόρθωση γωνιών και E-slice \cite{thistlethwaite1981,kociemba_twophase_details}.


\section{Φάση 2: Τετράδες και Ισοτιμία}
\label{sec:phase2}

\subsection{Περιορισμός Τετράδων Γωνιών}

Οι 8 γωνίες χωρίζονται σε δύο \textbf{τετράδες}:
\begin{itemize}
    \item \textbf{Τετράδα 1}: URF, UFL, DFR, DLF (μπροστινές γωνίες)
    \item \textbf{Τετράδα 2}: ULB, UBR, DBL, DRB (πίσω γωνίες)
\end{itemize}

Ο στόχος είναι οι γωνίες κάθε τετράδας να βρίσκονται σε θέσεις της ίδιας τετράδας (όχι απαραίτητα στη σωστή θέση) \cite{thistlethwaite1981}.

Αριθμός τρόπων τοποθέτησης:
\begin{equation}
    \binom{8}{4} = 70 \text{ καταστάσεις}
\end{equation}

\subsection{Διαχωρισμός Ακμών σε Slices}

Επιπλέον, οι ακμές πρέπει να παραμείνουν στα αντίστοιχα slices:
\begin{itemize}
    \item \textbf{UD-slice}: ακμές των εδρών U και D (θέσεις 0-7)
    \item \textbf{E-slice}: μεσαίες ακμές (θέσεις 8-11)
\end{itemize}

\subsection{Περιορισμός Ισοτιμίας}

Η \textbf{ισοτιμία μετάθεσης} (parity) πρέπει να είναι άρτια: η ισοτιμία γωνιών ισούται με την ισοτιμία ακμών. Οι στροφές 90° αλλάζουν την ισοτιμία, ενώ οι διπλές ($X2$) τη διατηρούν \cite{thistlethwaite1981,kociemba_twophase_details}.

\subsection{Περιορισμός Κινήσεων}

Αφαιρούνται οι $L$, $L'$, $R$, $R'$ (διατηρούνται $L2$, $R2$):
\begin{equation}
    \{U, U', U2, D, D', D2, F2, B2, L2, R2\} \quad (10 \text{ κινήσεις})
\end{equation}


\section{Φάση 3: Τελική Επίλυση}
\label{sec:phase3}

\subsection{Μόνο Διπλές Κινήσεις}

Στην τελική φάση επιτρέπονται μόνο οι 6 διπλές κινήσεις:
\begin{equation}
    \{U2, D2, F2, B2, L2, R2\}
\end{equation}

Αυτές οι κινήσεις διατηρούν \textbf{όλα} τα προηγούμενα αμετάβλητα \cite{thistlethwaite1981,kociemba_twophase_details}:
\begin{itemize}
    \item Προσανατολισμό ακμών (Φάση 0)
    \item Προσανατολισμό γωνιών (Φάση 1)
    \item E-slice τοποθέτηση (Φάση 1)
    \item Τετράδες και ισοτιμία (Φάση 2)
\end{itemize}

\subsection{Χώρος Μεταθέσεων}

Ο χώρος αναζήτησης περιορίζεται στις μεταθέσεις εντός κάθε τετράδας/slice:
\begin{itemize}
    \item Μεταθέσεις γωνιών εντός τετράδων: $(4!)^2 / 2 = 288$ (λόγω ισοτιμίας)
    \item Μεταθέσεις ακμών εντός slices: $(4!)^2 \times (4!)^2 / 2$
\end{itemize}

\subsection{Στόχος}

Η φάση ολοκληρώνεται όταν ο κύβος είναι πλήρως λυμένος:
\begin{equation}
    \text{cube.is\_solved()} = \text{True}
\end{equation}


\section{Κατασκευή Pattern Databases}
\label{sec:thistlethwaite_pdb}

\subsection{Μέθοδος BFS}

Στην παρούσα υλοποίηση, κάθε φάση χρησιμοποιεί \textbf{Pattern Database} (PDB)
που κατασκευάζεται με αναζήτηση κατά πλάτος (BFS) από τη λυμένη κατάσταση.
Η διαδικασία ακολουθεί τη γενική μεθοδολογία των pattern databases
\cite{culberson1998pattern}:

\begin{algorithm}[H]
\caption{Κατασκευή Pattern Database}
\begin{algorithmic}[1]
\State $\text{distance}[0] \gets 0$ \Comment{Λυμένη κατάσταση}
\State $\text{queue} \gets \{0\}$
\While{queue $\neq \emptyset$}
    \State $\text{coord} \gets$ queue.dequeue()
    \For{κάθε κίνηση $m$ στις επιτρεπόμενες}
        \State $\text{new\_coord} \gets$ apply($m$, coord)
        \If{distance[new\_coord] δεν έχει οριστεί}
            \State $\text{distance}[\text{new\_coord}] \gets \text{distance}[\text{coord}] + 1$
            \State queue.enqueue(new\_coord)
        \EndIf
    \EndFor
\EndWhile
\end{algorithmic}
\end{algorithm}

\subsection{Προδιαγραφές Βάσεων Δεδομένων}

\begin{table}[h]
\centering
\caption{Μεγέθη Pattern Databases ανά φάση}
\label{tab:pdb_sizes}
\begin{tabular}{clrrr}
\toprule
\textbf{Φάση} & \textbf{Συντεταγμένη} & \textbf{Μέγεθος} & \textbf{Μέγ. Βάθος} & \textbf{Μνήμη} \\
\midrule
0 & Προσ. ακμών & 2.048 & 7 & $\sim$2 KB \\
1 & Γωνίες $\times$ E-slice & 1.082.565 & 10 & $\sim$1 MB \\
2 & Τετράδες & 29.400 & 13 & $\sim$29 KB \\
3 & Μετάθεση γωνιών & 40.320 & 15 & $\sim$40 KB \\
\bottomrule
\end{tabular}
\end{table}

Τα PDBs αποθηκεύονται σε αρχεία για επαναχρησιμοποίηση. Η αρχική κατασκευή απαιτεί λίγα δευτερόλεπτα.


\section{Αναζήτηση IDA* ανά Φάση}
\label{sec:thistlethwaite_ida}

\subsection{Αλγόριθμος Αναζήτησης}

Κάθε φάση χρησιμοποιεί τον αλγόριθμο IDA* (Κεφάλαιο~\ref{ch:background}) με \cite{korf1985depth}:
\begin{itemize}
    \item \textbf{Ευρετική}: η τιμή από το αντίστοιχο Pattern Database
    \item \textbf{Επιτρεπόμενες κινήσεις}: το σύνολο κινήσεων της φάσης
    \item \textbf{Στόχος}: η συνθήκη μετάβασης στην επόμενη υποομάδα
\end{itemize}

\begin{algorithm}[H]
\caption{IDA* για φάση $i$}
\begin{algorithmic}[1]
\State $\text{bound} \gets \text{PDB}_i[\text{coord}]$
\Loop
    \State $\text{result} \gets$ SEARCH(cube, 0, bound, phase\_moves[$i$])
    \If{result = FOUND}
        \State \Return λύση φάσης
    \EndIf
    \State $\text{bound} \gets \text{result}$
\EndLoop
\end{algorithmic}
\end{algorithm}

\subsection{Αποκοπή Περιττών Κινήσεων}

Για βελτίωση απόδοσης, αποκλείονται περιττές ακολουθίες:
\begin{itemize}
    \item \textbf{Ίδια έδρα}: $U$ μετά $U'$ είναι περιττό (= $I$)
    \item \textbf{Αντίθετες έδρες}: επιβάλλεται σειρά $U$ πριν $D$ για κανονικοποίηση
    \item \textbf{Τριπλές επαναλήψεις}: $U U U = U'$
\end{itemize}

Αυτές οι αποκοπές μειώνουν σημαντικά τον παράγοντα διακλάδωσης.


\section{Σύνοψη και Αξιολόγηση}
\label{sec:thistlethwaite_summary}

\subsection{Χαρακτηριστικά Αλγορίθμου}

\begin{itemize}
    \item \textbf{Εγγύηση}: Λύση σε $\leq 52$ κινήσεις (τυπικά 30-45)
    \item \textbf{Ντετερμινισμός}: Πάντα βρίσκει λύση, χωρίς τυχαιότητα
    \item \textbf{Μνήμη}: Χαμηλές απαιτήσεις ($\sim$1.2 MB για όλα τα PDBs)
    \item \textbf{Χρόνος}: Μέτριος (τυπικά 1-5 δευτερόλεπτα)
\end{itemize}

\subsection{Πλεονεκτήματα και Μειονεκτήματα}

\textbf{Πλεονεκτήματα}:
\begin{itemize}
    \item Κομψή μαθηματική δομή βασισμένη στη θεωρία ομάδων
    \item Εκπαιδευτική αξία για κατανόηση της αλγεβρικής δομής του κύβου
    \item Αξιόπιστη και προβλέψιμη συμπεριφορά
\end{itemize}

\textbf{Μειονεκτήματα}:
\begin{itemize}
    \item Μακρύτερες λύσεις σε σχέση με σύγχρονους αλγορίθμους
    \item Η τετραφασική δομή δεν συνεπάγεται κατ' ανάγκην καλύτερο πρακτικό χρόνο από μια ώριμη two-phase υλοποίηση
    \item Τέσσερις φάσεις αντί δύο (περισσότερη επεξεργασία)
\end{itemize}

\subsection{Μετάβαση στον Kociemba}

Ο αλγόριθμος Kociemba (Κεφάλαιο~\ref{ch:kociemba}) βελτιώνει τον Thistlethwaite συμπτύσσοντας τις τέσσερις φάσεις σε δύο, επιτυγχάνοντας:
\begin{itemize}
    \item Λύσεις σημαντικά κοντύτερες από του Thistlethwaite και, γενικά, κοντά στο βέλτιστο
    \item Πιο συμπαγή two-phase αναζήτηση, με τα ακριβή measured timings του παρόντος artifact να αναφέρονται χωριστά στο Κεφάλαιο~\ref{ch:evaluation}
    \item Ευρεία χρήση σε εφαρμογές speedcubing
\end{itemize}

\chapter{Αλγόριθμος Kociemba}
\label{ch:kociemba}

% Reference: src/kociemba/, tests/unit/test_kociemba.py

\section{Περιγραφή του Αλγορίθμου}
% 2 phases, improvement over Thistlethwaite

\section{Σύστημα Συντεταγμένων}
% 6 coordinates, CoordCube

\section{Move Tables και Pruning Tables}
% Pre-computed tables, ~80MB

\section{Υλοποίηση και Αποτελέσματα}
% Two-phase IDA*, <19 moves

\chapter{Αλγόριθμος Korf - Βέλτιστη Επίλυση}
\label{ch:korf}

% Reference: src/korf/, tests/unit/test_a_star_solvers.py

\section{Εισαγωγή στη Βέλτιστη Επίλυση}
% God's Number = 20, why optimal is hard

\section{Pattern Databases για Βέλτιστη Επίλυση}
% Corner DB, Edge DBs, disjoint property

\section{IDA* με Pattern Database Heuristics}
% Additive heuristics, admissibility

\section{Αποτελέσματα}
% Optimal solutions ≤20, time analysis

\chapter{Εκτίμηση Απόστασης και Ευρετικές Συναρτήσεις}
\label{ch:heuristics}

% Reference: src/korf/heuristics.py, src/korf/composite_heuristic.py

\section{Το Πρόβλημα της Εκτίμησης Απόστασης}
% Definition, usefulness

\section{Υλοποιημένες Ευρετικές}
% Hamming, Manhattan, Pattern DB, Composite

\section{Σύνθετη Ευρετική (Composite Heuristic)}
% YOUR NOVEL CONTRIBUTION - expand this section!
% Dynamic strategy selection, entropy analysis

\section{Σύγκριση A* vs IDA*}
% Memory vs time tradeoff

\chapter{Πειραματική Αξιολόγηση}
\label{ch:evaluation}

% Reference: figures/, thesis_data_*.csv

\section{Μεθοδολογία}
% Test set, metrics, environment

\section{Σύγκριση Αλγορίθμων}
% Tables and figures from figures/

\section{Συζήτηση}
% When to use each algorithm

\chapter{Υλοποίηση}
\label{ch:implementation}

% Reference: entire src/, webapp/, ui/

\section{Αρχιτεκτονική Κώδικα}
% Project structure, modules

\section{Web Εφαρμογές}
% Next.js, Streamlit

\section{Testing}
% 203 tests, pytest

\chapter{Συμπεράσματα}
\label{ch:conclusions}

\section{Σύνοψη Αποτελεσμάτων}
% All objectives achieved

\section{Συνεισφορές}
% Implementation, composite heuristic, educational materials

\section{Μελλοντική Εργασία}
% Larger cubes, GPU, ML approaches


% ----------------------------------------------------------------------------
% BACK MATTER
% ----------------------------------------------------------------------------
\backmatter
\addcontentsline{toc}{chapter}{Βιβλιογραφία}
\bibliographystyle{IEEEtran}
\bibliography{references}

\appendix
\chapter{Οδηγός Εγκατάστασης και Χρήσης}
\label{app:usage}

\section{Απαιτήσεις Συστήματος}

Για την εκτέλεση του project απαιτείται:
\begin{itemize}
    \item Python 3.10 ή νεότερο
    \item \texttt{pip} για εγκατάσταση εξαρτήσεων
    \item βασικές βιβλιοθήκες επιστημονικού Python stack:
    \texttt{numpy}, \texttt{scipy}, \texttt{pandas}, \texttt{matplotlib}
    \item προαιρετικά: \texttt{streamlit} για το Python UI
    \item προαιρετικά: \texttt{node/npm} για τη Next.js εφαρμογή
\end{itemize}

Η τρέχουσα έκδοση του αποθετηρίου έχει επαληθευθεί με \texttt{Python 3.12.2}. Για πλήρη έλεγχο του περιβάλλοντος διατίθεται το script \filepath{verify_setup.py}.

\section{Εγκατάσταση}

\subsection{Βασική Εγκατάσταση}

\begin{lstlisting}
python3 -m venv venv
source venv/bin/activate
pip install -r requirements.txt
\end{lstlisting}

\subsection{Έλεγχος Περιβάλλοντος}

\begin{lstlisting}
python verify_setup.py
\end{lstlisting}

\subsection{Εκτέλεση Test Suite}

\begin{lstlisting}
python -m pytest tests -q
\end{lstlisting}

\section{Βασική Χρήση}

Η απλούστερη χρήση του project είναι μέσω των demo scripts. Τα scripts αυτά καλύπτουν δημιουργία κύβου, scrambling, επίλυση και σύγκριση αλγορίθμων.

\subsection{Εκτέλεση Demos}
\begin{lstlisting}
python demos/basic_usage.py
python demos/kociemba_demo.py
python demos/thistlethwaite_demo.py
python demos/a_star_comparison_demo.py
python demos/distance_estimator_demo.py
\end{lstlisting}

\subsection{Benchmarks και Συγκρίσεις}

\begin{lstlisting}
python demos/phase9/benchmark_demo.py --n-scrambles 20
python demos/phase9/algorithm_comparison_cli.py --depth 10 --seed 42
\end{lstlisting}

\section{Διαδραστικές Εφαρμογές}

\subsection{Streamlit}

\begin{lstlisting}
streamlit run ui/app.py
\end{lstlisting}

\subsection{Next.js Web Application}

\begin{lstlisting}
cd webapp
npm ci
npm run build
npm run dev
\end{lstlisting}

Η εφαρμογή Next.js λειτουργεί ως presentation/demo frontend και όχι ως authoritative live διεπαφή εκτέλεσης των Python solvers. Παρέχει σελίδες για:
\begin{itemize}
    \item αρχική παρουσίαση του project,
    \item demo επίλυσης μεμονωμένου scramble,
    \item demo σύγκρισης αλγορίθμων,
    \item και εκπαιδευτικό περιεχόμενο.
\end{itemize}
Η ζωντανή εκτέλεση των αλγορίθμων γίνεται από το \filepath{ui/} και από τα scripts benchmark/validation του repository.

\section{Αναπαραγωγή Αποτελεσμάτων Διπλωματικής}

Για την ακριβή αναπαραγωγή των αποτελεσμάτων του Κεφαλαίου~\ref{ch:evaluation}
χρησιμοποιείται ως επίσημη διαδρομή το script
\filepath{scripts/benchmarks/regenerate_thesis_benchmarks.py}, το οποίο
επαναεκτελεί τα benchmarks πάνω στο αποθηκευμένο σύνολο scrambles και
αναδημιουργεί τα αρχεία
\filepath{results/benchmarks/thesis/thesis_bench_d*.json} και
\filepath{results/benchmarks/thesis/thesis_results_combined.json}. Προαιρετικά
μπορούν να εκτελεστούν και τα scripts ανάλυσης, τα οποία διαβάζουν τα παραπάνω
αρχεία από τον ίδιο κατάλογο:

\begin{lstlisting}
python scripts/benchmarks/regenerate_thesis_benchmarks.py
python scripts/benchmarks/analyze_thesis_data.py
python scripts/benchmarks/generate_latex_tables.py
\end{lstlisting}

Τα scripts \filepath{generate_thesis_data.py} και
\filepath{generate_complete_thesis_data.py} διατηρούνται μόνο για ιστορικούς ή
διερευνητικούς λόγους και δεν αποτελούν τη διαδρομή αναπαραγωγής της
διπλωματικής.

Για τη διαχείριση της συγγραφής και των chapter packets προστέθηκε επίσης workflow script:

\begin{lstlisting}
python scripts/thesis_workflow.py status
python scripts/thesis_workflow.py packets --remaining
python scripts/thesis_workflow.py validate
\end{lstlisting}

\section{Σημειώσεις}

\begin{itemize}
    \item Η πρώτη εκτέλεση ορισμένων solvers μπορεί να είναι αισθητά πιο αργή λόγω δημιουργίας ή φόρτωσης πινάκων.
    \item Τα benchmark scripts αποθηκεύουν ρητά τα scrambles και τα αποτελέσματα σε JSON μορφή.
    \item Η πλήρης παραγωγή PDF της διπλωματικής απαιτεί TeX toolchain συμβατό με το \filepath{thesis/main.tex}. Στο αποθετήριο χρησιμοποιείται \texttt{tectonic}.
\end{itemize}

\chapter{Επιλεγμένος Κώδικας}
\label{app:code}

% Include key code snippets here
% Don't include entire files - just highlights

\section{Κλάση RubikCube}
% Simplified version from src/cube/rubik_cube.py

\section{Αλγόριθμος IDA*}
% From src/korf/a_star.py

\section{Σύνθετη Ευρετική}
% From src/korf/composite_heuristic.py


\end{document}
